\section{Finite examples and the scope of the size bounds}
\label{sec:costs}

The explicit family in Theorem~\ref{thm:quantitative} proves convergence
with every finite charge retained. Its constants are conservative, and
its guaranteed indices are far beyond useful sizes. Smaller fixed indices
can be evaluated directly from their rounded counts. This section gives certified finite examples and explains their
verification by exact-arithmetic enclosures.

\subsection{Bounds for evaluating finite types}

The entropy bounds in Section~\ref{sec:extraction} can be sharpened
without evaluating large factorials. Put \(c_0=(1+\log2)/2\).
For every integer \(t\ge1\),
\begin{equation}\label{eq:elementary-factorial}
 (t+\tfrac12)\log t-t+c_0
 \le\log(t!)\le(t+\tfrac12)\log t-t+1.
\end{equation}
For the upper bound, the trapezoidal sum of the concave function
\(\log x\) on \([1,t]\) is at most its integral. For the lower
bound, Jensen's inequality on each interval
\([i-1/2,i+1/2]\) bounds its integral by \(\log i\).
Sum over \(i=1,\ldots,t\), evaluate the integral, and use
\((t+1/2)\log(1+1/(2t))\ge1/2\), a consequence of
\(\log(1+x)\ge x/(1+x)\). This proves both inequalities.

For an action with \(s\) positive child counts, let
\[
 C_i=\frac12\log(Nw_i)
       -\frac12\sum_{y:n_{i,y}>0}\log(Nn_{i,y}).
\]
Subtracting the factorial bounds gives
\begin{equation}\label{eq:finite-multinomial}
 NH_i+C_i+c_0-s\le\log M_i
       \le NH_i+C_i+1-sc_0.
\end{equation}
Intersect these bounds with Lemma~\ref{lem:types} and with
\(\log M_i\ge0\). A type with one positive child count has
\(\log M_i=0\) exactly. For a reset apply
\eqref{eq:finite-multinomial} to the two counts
\((m_r,Nz_r-m_r)\), omitting any zero count, then add \((Nz_r-m_r)\log2\).
Intersect the result with
\([Nz_r\log3-\log(Nz_r+1),\,Nz_r\log3]\).
This uses the actual modal integer \(m_r\), without replacing it
by \(Nz_r/3\).

These sector bounds give a lower bound for \(\Delta_N\) and an
upper bound for \(\log E\). Together with the actual signed
catalyst in \eqref{eq:kappa} and
\(\log(S/H)<3I\log2\), they give a sufficient upper bound for
\(\Phi\). All expressions involve logarithms of integer counts and
small rank bases. They do not require the factorials, sector words, or
coefficient arrays to be materialized.

\subsection{Saved finite certificates}

Table~\ref{tab:finite-multiway} records three fixed comparisons at the
larger rounding scale
\(K=a^3\lceil\log_2a\rceil^2\), \(N=1\), \(k=32\), and
\(j=128a\lceil\log_2a\rceil\). Increasing \(j\) amortizes
catalyst supply; increasing \(k\) amortizes interpolation. Both also
increase the source dimension \(A=u^{kj}\).

\begin{table}[ht]
\centering
\begin{tabular}{rrrrr}
\toprule
\(a\)&\(K\)&\(j\)&upper bound for \(\Phi\)&decimal digits of \(A\)\\
\midrule
4&256&1\,024&3.628926097&84\,456\,434\\
8&4\,608&3\,072&2.848599165&11\,877\,489\,885\\
12&27\,648&6\,144&2.675820359&234\,905\,651\,862\\
\bottomrule
\end{tabular}
\caption{Certified sufficient exponent bounds at fixed small indices.
Here \(N=1\) and \(k=32\). Each exponent is rounded upwards from
its rational interval. The digit counts refer to the recursive source
side, not to a stored or executed matrix.}
\label{tab:finite-multiway}
\end{table}

The last row proves that this finite relative scheme has pooled exponent
less than \(2.675820359\). Its size also shows why this is a mathematical
construction rather than an implementation result. The table does not
bound the exact exponents from below, compare practical running times,
or give the dimensions of the final rational seeds.

The saved certificates use rational logarithm enclosures. For a positive
integer write it as \(2^e y\), \(1\le y\le2\), and bracket
\(y\) by adjacent rationals of denominator \(2^{80}\). For
\(z=(y-1)/(y+1)\),
\begin{equation}\label{eq:log-enclosure}
 2\sum_{h=0}^{39}\frac{z^{2h+1}}{2h+1}
 \le\log y\le
 2\sum_{h=0}^{39}\frac{z^{2h+1}}{2h+1}
 +\frac{2z^{81}}{81(1-z^2)}.
\end{equation}
The positive tail of the power series gives the remainder bound.
Evaluate at the rational bracket endpoints, apply the same bound to
\(\log2\), and combine the intervals with outward rounding.
For a positive rational, subtract the denominator's logarithm interval
from the numerator's. Digit counts follow by enclosing \(\log_{10}A\)
strictly between consecutive integers.
The supplementary report \texttt{results/multiway-finite-costs.json}
contains the integer types, signed catalyst counts, charge intervals,
and source hashes underlying the table.

\subsection{Comparison with exact normalization}

Appendix~\ref{sec:exact-certificate} gives a closed binary certificate
and its exact minimal denominator \(K_a\). It satisfies
\[
 K_a\ge\frac32\lcm(1,\ldots,a-1)\ge\frac{3\,2^{a-2}}a.
\]
This is an obstruction to using those exact rational frequencies at a
smaller scale. Conserved rounding avoids that requirement: its edge
counts need only approximate the frequencies while preserving the entire
resource identity. The bound \(\mathcal L\le4J_a\) measures the
resulting loss. Multiway relations also reduce the ideal total auxiliary
charge to \(O(a)\). Together these changes give the bound
\(\log A=O(a^5(\log a)^2)\).

Threefold balancing separately reduces the dimension charges relative to
sixfold balancing at the same \(k\). It does not by itself halve an
exponent: doubling \(k\) in the threefold construction recovers the
same charged scalar counts as using \(k\) with sixfold balancing,
when the underlying types agree. Likewise, increasing a sector count
improves the target dimension but also increases the separator degree.
The full expression \eqref{eq:pooled-exponent} is needed to assess either
change.

The dimension \(A\) belongs to the specified relative scheme. It is
not a lower bound on all algorithms obtainable by restricting its source,
or on all rational seeds with the same exponent. For example, tensoring
a small scheme with \(\MM_C\) increases its displayed source side,
while restricting that common factor to one coefficient recovers the
small scheme. A source-size obstruction therefore needs a stated
construction and reduction model. Neither the denominator benchmark nor
the displayed finite examples imply a general size barrier for fast
matrix multiplication.
